\begin{center}
    \large
    \begin{singlespace}    
        \textbf{\chineseTitle{}} \\[0.5cm]
    \end{singlespace}

    \begin{singlespace}    
    學生：\studentCnName  \hspace{2.5cm}  指導教授：\advisorCnName \hspace{0.1cm} 博士 \\

    \ifdefined\advisorCnNameB % 如果有共同指導教授
        \hspace{9.6cm}  \advisorCnNameB ~ 博士 \\
    \fi
    \end{singlespace}

    \vspace{0.5cm}

    \begin{singlespace}
        國立陽明交通大學\\
        \NameofDepartmentInstituteCN\\[0.2cm]
    \end{singlespace}

    \textbf{摘\hspace{1cm}要} \\[0.5cm]

\end{center}

\normalsize 

過渡金屬硫屬化物 (Transition metal dichalcogenides, TMDs) 的能谷激子天生帶有自旋角動量 (SAM)，而軌道角動量 (OAM) 則能提供理論上無限維度的調控自由度；因此，將能谷 SAM 高效且具選擇性地轉化為 OAM，是發展高容量量子光源與谷光子學應用的新興課題。

本研究提出一套反向設計框架，結合三維時域差分 (Finite-Difference Time-Domain, FDTD) 方法與伴隨變量法 (Adjoint Method)，設計出高自由度的磷化鎵 ($\text{GaP}$) 拓撲最佳化 (Topology Optimization) 結構。此結構能在固定的極小體積內，將單層二硫化鎢 ($\text{WS}_2$) 激子的能谷 SAM 確定性地轉換為任意指定拓撲荷的 OAM，產生高純度的拉格耳-高斯（Laguerre-Gaussian, LG）光束。

數值模擬以 $\ell = 0 \text{--} 5$ 的完整系列進行證明：在 $\text{WS}_2$ 激子發光波長下，針對左旋圓偏振 (LCP) 之發光進行結構優化，最終元件的遠場模態純度皆可達到 $80\%$ 以上。此外，當激子發光轉變為右旋圓偏振 (RCP) 時，目標 LG 模態的遠場純度降至 $1\%$ 以下，展現出極高的谷極化選擇性。此成果為微型化渦旋光學元件與谷電子學元件的設計提供了全新方案。
\vspace{1cm}
% 中文摘要及關鍵詞 5-7 個 

\textbf{關鍵字：}谷極化, 過渡金屬硫屬化物, 拉格耳高斯光束, 超表面, 拓樸最佳化, 伴隨變量法